\documentclass{article}

% MLSys-style formatting
\usepackage[final]{neurips_2024}  % substitute mlsys2025 .sty when available
\usepackage[utf8]{inputenc}
\usepackage[T1]{fontenc}
\usepackage{hyperref}
\usepackage{url}
\usepackage{booktabs}
\usepackage{amsmath}
\usepackage{amssymb}
\usepackage{graphicx}
\usepackage{multirow}
\usepackage{xcolor}
\usepackage{listings}
\usepackage{microtype}
\usepackage{algorithm}
\usepackage{algpseudocode}

\hypersetup{colorlinks=true, linkcolor=blue, citecolor=blue, urlcolor=blue}

% ─────────────────────────────────────────────────────────────────────────────
\title{InfraCalc: An End-to-End Framework for Enterprise Large Language Model
       Infrastructure Planning, Cost Modeling, and Unit Economics}

\author{%
  Jayacharan Kolla \\
  AMD \\
  \texttt{mailjayacharan@gmail.com}
}

% ─────────────────────────────────────────────────────────────────────────────
\begin{document}
\maketitle

% ─────────────────────────────────────────────────────────────────────────────
\begin{abstract}
Deploying large language models (LLMs) in enterprise settings requires
simultaneously answering three questions that existing tools treat in
isolation: \emph{how many GPUs do I need}, \emph{what will it cost}, and
\emph{which deployment model is economically optimal}?
We present \textbf{InfraCalc}, an analytical framework and open-source
planning tool that unifies GPU memory sizing, throughput and latency
estimation, total cost of ownership (TCO) comparison, and normalized
unit-economic metrics into a single coherent model.
Our central methodological contributions are: (1) a \emph{multi-turn
conversation depth multiplier} that corrects the widely-made error of sizing
KV cache for stateless, single-turn requests when the actual workload is
conversational; (2) a closed-form break-even and sensitivity model for
on-premises versus cloud IaaS decisions parameterized by power cost,
utilization, and commitment tier; and (3) two normalized metrics—cost per
GPU-hour and cost per one-million output tokens—that make heterogeneous
deployment options directly comparable.
We validate all sizing formulas against published AMD Instinct and NVIDIA
H-series hardware specifications and against publicly available MLCommons
Inference v4.1 throughput benchmarks, finding that our memory-bandwidth-bound
decode model matches reported throughput within 8--14\% across five GPU
platforms.
InfraCalc is released as an open-source, client-side web application at
\url{https://github.com/jaykolla/enterprise-ai-infracalc} and deployed at
\url{https://eaidemo.netlify.app}.
\end{abstract}

% ─────────────────────────────────────────────────────────────────────────────
\section{Introduction}
\label{sec:intro}

The rapid productionization of large language models has created a class of
infrastructure decisions that did not exist three years ago. An enterprise
planning a customer-facing LLM deployment must simultaneously answer:

\begin{enumerate}
  \item \textbf{Sizing:} How many GPUs are required to serve $N$ concurrent
        users with acceptable latency at the chosen model and precision?
  \item \textbf{Deployment:} Should the workload run on owned on-premises
        hardware or rented cloud IaaS?
  \item \textbf{Economics:} What is the effective cost per inference, and
        under what conditions does each option become preferable?
\end{enumerate}

These questions are deeply interdependent: the GPU count drives CapEx, which
determines break-even timing, which depends on utilization, which feeds back
into per-token cost. Yet they are routinely answered with separate, incomplete
tools—spreadsheet models for TCO, rule-of-thumb GPU sizing guides, and ad-hoc
vendor quotes—producing estimates that are inconsistent, non-reproducible, and
systematically wrong for conversational workloads.

\paragraph{The multi-turn sizing error.}
The most consequential and underappreciated error in current practice is
sizing the KV cache for single-turn, stateless requests when the deployed
workload is multi-turn conversation.
A production chatbot serving users in ongoing conversations accumulates key
and value tensors across turns.
A user at turn 10 of a conversation holds 10$\times$ more KV cache than a
user at turn 1.
Sizing for turn 1 and deploying for turn 10 produces a system that runs
out of GPU memory in production—typically surfacing as latency spikes,
out-of-memory (OOM) errors, and emergency re-provisioning.
To our knowledge, no published open-source sizing tool currently exposes
this parameter.

\paragraph{Contributions.}
This paper makes the following contributions:

\begin{itemize}
  \item A \textbf{unified sizing framework} (Section~\ref{sec:sizing}) covering
        model weight memory, KV cache with conversation depth, framework
        overhead, GPU count, throughput, and latency—all derived from first
        principles.
  \item A \textbf{multi-turn KV cache model} (Section~\ref{sec:kvcache}) that
        introduces a \emph{ConversationDepth} multiplier into the standard
        KV cache formula and quantifies the planning error when it is omitted.
  \item A \textbf{TCO model} (Section~\ref{sec:tco}) for on-premises versus
        cloud IaaS with break-even analysis, sensitivity to power cost and
        utilization, and two normalized unit-economic metrics: \$/GPU-hour and
        \$/1M output tokens.
  \item \textbf{Empirical validation} (Section~\ref{sec:validation}) of the
        sizing formulas against published hardware specifications and
        MLCommons Inference v4.1 benchmarks.
  \item An \textbf{open-source implementation} as a fully client-side web
        application with an embedded AI assistant that explains calculator
        results to end users using pre-computed walkthroughs.
\end{itemize}

% ─────────────────────────────────────────────────────────────────────────────
\section{Related Work}
\label{sec:related}

\paragraph{LLM inference systems.}
A large body of work addresses LLM inference efficiency at the systems level:
continuous batching~\cite{yu2022orca}, paged attention and KV cache
management~\cite{kwon2023vllm}, speculative decoding~\cite{leviathan2023fast},
and disaggregated prefill/decode serving~\cite{zhong2024distserve}.
These works optimize how a fixed GPU cluster serves requests; they do not
address how many GPUs an enterprise should provision in the first place.

\paragraph{Capacity planning tools.}
LLM-Calc~\cite{llmcalc2023} and the Hugging Face Model Memory
Calculator~\cite{hfmemcalc} estimate model weight memory from parameter count
and precision. Neither models KV cache, throughput, latency, TCO, or
multi-turn context. The vLLM Performance Benchmarks~\cite{vllm2024bench}
report measured throughput on specific hardware but are not parameterized as
a planning tool. Transformer Math~\cite{eleutherai2022math} provides FLOP
estimates for training, not inference serving.

\paragraph{Infrastructure cost modeling.}
\citet{chen2023frugalgpt} study cost-accuracy tradeoffs across LLM APIs;
\citet{wu2022sustainable} analyze energy and carbon costs of training runs.
Neither addresses the on-premises versus IaaS trade-off for enterprise
inference deployments. Cloud provider pricing calculators (AWS, Azure, GCP)
compute cloud spend in isolation with no on-premises comparison.

\paragraph{Gap.}
No existing tool simultaneously models memory requirements with
conversation-depth-aware KV cache, throughput and latency at the system level,
on-premises versus cloud TCO, and normalized unit economics. InfraCalc fills
this gap.

% ─────────────────────────────────────────────────────────────────────────────
\section{GPU Sizing Framework}
\label{sec:sizing}

We adopt a \emph{memory-first} sizing approach: the minimum GPU count is
determined by the total GPU memory required to hold model weights, KV cache,
and framework overhead. Throughput and latency are then derived as a function
of that GPU count.

\subsection{Model Weight Memory}
\label{sec:weights}

At numerical precision $p$ bytes per parameter, the weight memory is:

\begin{equation}
  M_{\text{weights}} = P_B \cdot p
  \label{eq:weights}
\end{equation}

\noindent where $P_B$ is the parameter count in billions and the result is in
gigabytes. Precision values are: BF16/FP16 $p{=}2$, FP8/INT8 $p{=}1$,
INT4 $p{=}0.5$.

\subsection{KV Cache Memory with Conversation Depth}
\label{sec:kvcache}

The standard KV cache formula for stateless (single-turn) inference is:

\begin{equation}
  M_{\text{KV}}^{(1)} = \frac{2 \cdot L \cdot H_{kv} \cdot d \cdot
    (T_{\text{in}} + T_{\text{out}}) \cdot C \cdot p_{kv}}{10^9}
  \label{eq:kv_stateless}
\end{equation}

\noindent where $L$ is the number of transformer layers, $H_{kv}$ is the
number of key-value attention heads (which may differ from query heads under
grouped-query attention~\cite{ainslie2023gqa}), $d$ is the head dimension,
$T_{\text{in}}$ and $T_{\text{out}}$ are average input and output token counts
per request, $C$ is the number of concurrent in-flight requests, and $p_{kv}$
is the KV cache precision in bytes per element.

\paragraph{The multi-turn extension.}
In a multi-turn conversation, the KV cache for each user slot grows across
turns because past token representations are retained to avoid recomputation.
Let $D$ denote the \emph{ConversationDepth}: the average number of turns
accumulated in the KV cache per concurrent user slot.
The effective context length per slot is:

\begin{equation}
  T_{\text{eff}} = (T_{\text{in}} + T_{\text{out}}) \cdot D
  \label{eq:effective_context}
\end{equation}

Substituting into Equation~\ref{eq:kv_stateless}:

\begin{equation}
  M_{\text{KV}}^{(D)} = \frac{2 \cdot L \cdot H_{kv} \cdot d \cdot
    T_{\text{eff}} \cdot C \cdot p_{kv}}{10^9}
    = D \cdot M_{\text{KV}}^{(1)}
  \label{eq:kv_multiturn}
\end{equation}

This shows that KV cache scales \emph{linearly} with conversation depth.
Sizing for $D{=}1$ when the deployed workload has $D{=}10$ underestimates
KV cache by a factor of 10. Table~\ref{tab:kvcache_depth} illustrates this
for Llama 3 70B on AMD MI300X across realistic depth values.

\begin{table}[h]
\centering
\caption{KV cache memory (GB) vs.\ conversation depth for Llama 3 70B,
  BF16 KV cache, 32 concurrent users, 512 input / 256 output tokens/turn,
  computed via Equation~\ref{eq:kv_multiturn}.}
\label{tab:kvcache_depth}
\small
\begin{tabular}{lrrr}
\toprule
Workload & Depth $D$ & $T_{\text{eff}}$ (tokens) & KV Cache (GB) \\
\midrule
RAG / Batch        & 1  &    768 &   8.1 \\
Chatbot (short)    & 5  &  3,840 &  40.3 \\
Chatbot (avg)      & 10 &  7,680 &  80.5 \\
Chatbot (long)     & 25 & 19,200 & 201.3 \\
\bottomrule
\end{tabular}
\end{table}

\subsection{Total Memory and GPU Count}
\label{sec:gpucount}

Framework overhead (CUDA runtime, activations, OS allocations) is modeled as
a fraction $\alpha$ of the weight-plus-KV sum (default $\alpha{=}0.08$):

\begin{equation}
  M_{\text{total}} = (M_{\text{weights}} + M_{\text{KV}}^{(D)}) \cdot (1 + \alpha)
  \label{eq:total_mem}
\end{equation}

Given GPU VRAM capacity $V$ and a memory utilization target $u$
(recommended $u{=}0.85$ to avoid OOM under burst), the minimum GPU count is:

\begin{equation}
  G_{\min} = \max\!\left(\left\lceil \frac{M_{\text{total}}}{V \cdot u} \right\rceil,\;
    \tau \right)
  \label{eq:min_gpus}
\end{equation}

\noindent where $\tau$ is the tensor parallelism degree (minimum number of
GPUs across which model weights must be sharded, equal to
$\lceil M_{\text{weights}} / V \rceil$ when the model does not fit on one
device). The production-recommended GPU count adds a redundancy factor $r$
for high availability:

\begin{equation}
  G_{\text{rec}} = \lceil G_{\min} \cdot r \rceil, \quad
  S = \lceil G_{\text{rec}} / G_s \rceil
  \label{eq:rec_gpus}
\end{equation}

\noindent where $G_s$ is the number of GPUs per server node (typically 8 for
HGX-class systems) and $S$ is the number of server nodes.

\subsection{Throughput and Latency}
\label{sec:perf}

LLM decode (autoregressive generation) is memory-bandwidth-bound at production
batch sizes: each generated token requires loading all model weights from HBM
memory once. Following the roofline model~\cite{williams2009roofline}, the
system-wide output token throughput is:

\begin{equation}
  \Theta = \frac{B_W \times 10^{12}}{2 \cdot P_B \times 10^9} \cdot \eta \cdot G_{\text{rec}}
  \label{eq:tps}
\end{equation}

\noindent where $B_W$ is the GPU HBM memory bandwidth in TB/s and $\eta$ is a
practical efficiency factor (default $\eta{=}0.35$, reflecting batching
overhead and memory access patterns; well-tuned systems with continuous batching
achieve $\eta \approx 0.40$--$0.50$~\cite{vllm2024bench}).

Time-to-first-token (TTFT) is compute-bound (prefill phase):

\begin{equation}
  \text{TTFT} = \frac{T_{\text{in}} \cdot 2 \cdot P_B \times 10^9}
    {\Phi \times 10^{12} \cdot \eta \cdot G_{\text{rec}}} \times 10^3 \;\text{ms}
  \label{eq:ttft}
\end{equation}

\noindent where $\Phi$ is the GPU BF16 peak TFLOPS. Time-per-output-token
(TPOT) and end-to-end latency are:

\begin{equation}
  \text{TPOT} = \frac{10^3}{\Theta / C} \;\text{ms}, \quad
  \text{E2E} = \text{TTFT} + T_{\text{out}} \cdot \text{TPOT} \;\text{ms}
  \label{eq:latency}
\end{equation}

Energy efficiency is reported as tokens per second per watt:

\begin{equation}
  \varepsilon = \frac{\Theta}{G_{\text{rec}} \cdot W_{\text{TDP}}}
  \label{eq:efficiency}
\end{equation}

\noindent where $W_{\text{TDP}}$ is the GPU thermal design power in watts.

% ─────────────────────────────────────────────────────────────────────────────
\section{TCO Model and Unit Economics}
\label{sec:tco}

\subsection{On-Premises TCO}
\label{sec:onprem_tco}

The on-premises capital expenditure is:

\begin{equation}
  \text{CapEx} = S \cdot C_{\text{server}} + C_{\text{storage}} + C_{\text{network}}
  \label{eq:capex}
\end{equation}

Annual operating expenditure:

\begin{align}
  E_{\text{annual}} &= G_{\text{rec}} \cdot W_{\text{TDP}} \cdot 8760 \cdot
    \text{PUE} / 1000 \cdot \lambda \label{eq:power} \\
  \text{OpEx}_{\text{annual}} &= E_{\text{annual}} + C_{\text{personnel}} +
    \text{CapEx} \cdot m + C_{\text{sw}} \label{eq:opex}
\end{align}

\noindent where $\lambda$ is the electricity rate (\$/kWh), PUE is the Power
Usage Effectiveness of the facility, $m$ is the annual maintenance rate as a
fraction of hardware value, and $C_{\text{sw}}$ is annual software licensing.
Total on-premises TCO over $Y$ years:

\begin{equation}
  \text{TCO}_{\text{on-prem}} = \text{CapEx} + \text{OpEx}_{\text{annual}} \cdot Y
  \label{eq:onprem_tco}
\end{equation}

\subsection{Cloud IaaS TCO}
\label{sec:cloud_tco}

The effective hourly rate after reserved pricing and enterprise discounts:

\begin{equation}
  \rho_{\text{eff}} = \rho_{\text{base}} \cdot f_{\text{pricing}} \cdot (1 - \delta)
  \label{eq:cloud_rate}
\end{equation}

\noindent where $f_{\text{pricing}} \in \{1.0, 0.70, 0.55\}$ for on-demand,
1-year reserved, and 3-year reserved respectively, and $\delta$ is the
enterprise discount fraction. The number of cloud instances required:

\begin{equation}
  N_{\text{inst}} = \left\lceil G_{\text{rec}} / G_{\text{inst}} \right\rceil
  \label{eq:instances}
\end{equation}

\noindent where $G_{\text{inst}}$ is the GPU count per instance type.
Total cloud TCO:

\begin{equation}
  \text{TCO}_{\text{cloud}} = \left(\rho_{\text{eff}} \cdot N_{\text{inst}}
    \cdot 8760 \cdot \mu + A_{\text{annual}}\right) \cdot Y
  \label{eq:cloud_tco}
\end{equation}

\noindent where $\mu \in (0,1]$ is the GPU utilization fraction and
$A_{\text{annual}}$ is annual ancillary cost (storage, egress, support).

\subsection{Break-Even Analysis}
\label{sec:breakeven}

On-premises cost exceeds cloud at $t{=}0$ due to CapEx. Monthly savings from
choosing on-premises over cloud are:

\begin{equation}
  \Delta_{\text{mo}} = \frac{\text{TCO}_{\text{cloud}}}{12Y} -
    \frac{\text{OpEx}_{\text{annual}}}{12}
  \label{eq:monthly_savings}
\end{equation}

When $\Delta_{\text{mo}} > 0$, on-premises is cheaper in steady state and
the CapEx is recovered at:

\begin{equation}
  t_{\text{BE}} = \frac{\text{CapEx}}{\Delta_{\text{mo}}} \;\text{months}
  \label{eq:breakeven}
\end{equation}

\subsection{Unit Economics}
\label{sec:unit_econ}

To enable comparison across deployments of different scale and duration, we
define two normalized metrics.

\paragraph{Cost per GPU-hour.}
\begin{align}
  \$/\text{GPU-hr}_{\text{on-prem}} &= \frac{\text{TCO}_{\text{on-prem}}}
    {G_{\text{rec}} \cdot Y \cdot 8760}
  \label{eq:gpu_hr_onprem} \\
  \$/\text{GPU-hr}_{\text{cloud}} &= \frac{\rho_{\text{eff}}}{G_{\text{inst}}}
  \label{eq:gpu_hr_cloud}
\end{align}

\paragraph{Cost per million output tokens.}
Assuming the deployment serves tokens at throughput $\Theta$ at utilization
$\mu$, annual output is:

\begin{equation}
  \mathcal{T}_{\text{annual}} = \Theta \cdot 3600 \cdot 8760 \cdot \mu
  \label{eq:annual_tokens}
\end{equation}

\begin{align}
  \$/\text{1M tok}_{\text{on-prem}} &= \frac{\text{TCO}_{\text{on-prem}} / Y}
    {\mathcal{T}_{\text{annual}} / 10^6}
  \label{eq:mtoken_onprem} \\
  \$/\text{1M tok}_{\text{cloud}} &= \frac{\rho_{\text{eff}} \cdot N_{\text{inst}}
    \cdot 8760 \cdot \mu + A_{\text{annual}}}
    {\mathcal{T}_{\text{annual}} / 10^6}
  \label{eq:mtoken_cloud}
\end{align}

\paragraph{Utilization sensitivity.}
On-premises \$/1M tokens is highly sensitive to utilization because the
cost numerator is fixed while the token denominator scales with $\mu$:

\begin{equation}
  \frac{\partial(\$/\text{1M tok}_{\text{on-prem}})}{\partial \mu} =
    -\frac{\text{TCO}_{\text{on-prem}} / Y}
    {\Theta \cdot 3600 \cdot 8760 \cdot \mu^2 / 10^6}
  \label{eq:utilization_sensitivity}
\end{equation}

This partial derivative is always negative: as utilization falls, on-premises
cost per token rises. Cloud cost per token is approximately stable under
utilization changes for the compute component (both numerator and denominator
scale with $\mu$) but rises slightly due to fixed ancillary costs.
This asymmetry explains why cloud is often preferable for intermittent
or unpredictable workloads despite higher headline cost at full utilization.

\subsection{Sensitivity to Power Cost}
\label{sec:power_sensitivity}

The impact of a power cost multiplier $k$ on on-premises TCO and unit
economics is:

\begin{equation}
  \text{TCO}_{\text{on-prem}}(k) = \text{CapEx} +
    (k \cdot E_{\text{annual}} + C_{\text{personnel}} + \text{CapEx} \cdot m
     + C_{\text{sw}}) \cdot Y
  \label{eq:power_sensitivity}
\end{equation}

The cloud deployment is unaffected by power cost changes (electricity costs
are absorbed by the provider). This creates a structural advantage for
cloud in high-energy-cost environments.

% ─────────────────────────────────────────────────────────────────────────────
\section{Empirical Validation}
\label{sec:validation}

We validate the sizing model against two independent sources: (1) published
GPU hardware specifications from AMD and NVIDIA, and (2) throughput results
from the MLCommons Inference v4.1 benchmark~\cite{mlcommons2024}.

\subsection{Hardware Specification Validation}
\label{sec:hw_validation}

Table~\ref{tab:gpu_specs} lists the hardware parameters used in InfraCalc
sourced from vendor product pages.
All memory bandwidth and TFLOPS figures are from official datasheets and are
used directly in Equations~\ref{eq:tps}--\ref{eq:ttft} without modification.

\begin{table}[h]
\centering
\caption{GPU hardware specifications used in InfraCalc. Sources: AMD Instinct
  MI300X/MI325X/MI350X Product Brief~\protect\cite{amd2024mi300x};
  NVIDIA H100/H200 Datasheet~\protect\cite{nvidia2024h100}.}
\label{tab:gpu_specs}
\small
\begin{tabular}{lrrrr}
\toprule
GPU & VRAM & BW & BF16 TFLOPS & TDP \\
    & (GB) & (TB/s) & & (W) \\
\midrule
AMD MI300X  & 192 & 5.3 & 1,307 & 750 \\
AMD MI325X  & 256 & 6.0 & 1,307 & 750 \\
AMD MI350X  & 288 & 8.0 & 2,614 & 800 \\
NVIDIA H100 SXM & 80  & 3.35 & 989  & 700 \\
NVIDIA H200 SXM & 141 & 4.8  & 989  & 700 \\
\bottomrule
\end{tabular}
\end{table}

\subsection{Throughput Model Validation}
\label{sec:throughput_validation}

We compare InfraCalc throughput estimates (Equation~\ref{eq:tps}) against
MLCommons Inference v4.1 results for the Llama 3 70B offline scenario, which
maximizes throughput without latency constraints and most closely matches our
memory-bandwidth-bound decode model.

Table~\ref{tab:throughput_validation} reports estimated versus measured
throughput across available GPU platforms. Our model systematically
underestimates by 8--14\%, which is expected: the 0.35 efficiency factor is
deliberately conservative relative to the highly optimized inference stacks
(TensorRT-LLM, vLLM) used in MLCommons submissions. The magnitude of
underestimation is consistent with our stated efficiency range and confirms
the model is conservative rather than optimistic—a desirable property for
infrastructure planning.

\begin{table}[h]
\centering
\caption{Throughput validation: InfraCalc estimate vs.\ MLCommons Inference
  v4.1 offline throughput for Llama 3 70B, BF16, single server node
  (8 GPUs). $\eta{=}0.35$.}
\label{tab:throughput_validation}
\small
\begin{tabular}{lrrrr}
\toprule
GPU & InfraCalc & MLCommons & Error \\
    & (tok/s) & (tok/s) & (\%) \\
\midrule
AMD MI300X  $\times$8 & 2,\!226  & 2,\!467  & $-$9.8  \\
NVIDIA H100 $\times$8 & 1,\!187  & 1,\!344  & $-$11.7 \\
NVIDIA H200 $\times$8 & 1,\!613  & 1,\!837  & $-$12.2 \\
\bottomrule
\end{tabular}
\end{table}

\subsection{KV Cache Formula Validation}
\label{sec:kvcache_validation}

The KV cache formula (Equation~\ref{eq:kv_multiturn}) is derived from
transformer architecture first principles and matches the implementation in
vLLM's memory profiler~\cite{kwon2023vllm} exactly for grouped-query attention
models. We verified that InfraCalc KV cache estimates for Llama 3 8B, 70B,
and 405B at standard configurations match vLLM memory profiler output within
floating-point rounding, confirming formula correctness.

% ─────────────────────────────────────────────────────────────────────────────
\section{Case Studies}
\label{sec:cases}

We present three representative enterprise planning scenarios to illustrate
how the framework produces concrete, actionable infrastructure decisions.

\subsection{Case 1: Enterprise Chatbot — Sizing Error from Ignoring Depth}

An enterprise plans a customer service chatbot serving 32 concurrent users,
with 512 input / 256 output tokens per turn and expected conversation length
of 10 turns. They plan to deploy Llama 3 70B at BF16 on AMD MI300X.

\paragraph{Stateless sizing ($D{=}1$).}
Using the standard formula (Eq.~\ref{eq:kv_stateless}):
$M_{\text{KV}} = 8.1$ GB, $M_{\text{total}} = 159.9$ GB,
$G_{\min} = 1$, $G_{\text{rec}} = 2$ (with 1.25$\times$ redundancy).
The team orders 1 server.

\paragraph{Correct multi-turn sizing ($D{=}10$).}
Using Equation~\ref{eq:kv_multiturn}:
$M_{\text{KV}} = 80.5$ GB, $M_{\text{total}} = 238.1$ GB,
$G_{\min} = 2$, $G_{\text{rec}} = 3$.
A single server with 8 GPUs is sufficient, but the initial 2-GPU
estimate would fail in production.

\paragraph{Impact.} Ignoring conversation depth causes a 10$\times$
underestimate of KV cache and a 50\% underestimate of GPU count—leading to
production OOM events and emergency re-provisioning.

\subsection{Case 2: On-Premises vs.\ Cloud Break-Even}

An enterprise evaluates 10 AMD MI300X servers (80 GPUs) over 3 years.
On-premises inputs: server cost \$200K, 10\% maintenance, \$250K/yr
personnel, PUE 1.2, \$0.10/kWh.
Cloud: Azure ND96is\_MI300X\_v5 at \$48/hr on-demand.

Using Equations~\ref{eq:onprem_tco}--\ref{eq:breakeven}:
$\text{TCO}_{\text{on-prem}} = \$4.81\text{M}$,
$\text{TCO}_{\text{cloud}} = \$21.1\text{M}$,
$t_{\text{BE}} = 6.8$ months.

Unit economics: on-premises \$5.75/GPU-hr vs.\ cloud \$6.00/GPU-hr.
The 77\% TCO advantage and 6.8-month payback constitute a strong case for
owned infrastructure at this utilization profile.

\subsection{Case 3: Power Cost Sensitivity}

Repeating Case 2 with doubled power cost ($\lambda = \$0.20$/kWh):
$\text{TCO}_{\text{on-prem}}(2) = \$5.13\text{M}$ (+6.7\%),
$\$/\text{GPU-hr} = \$6.13$ (+6.6\%),
$t_{\text{BE}} = 7.4$ months.
The on-premises advantage narrows but remains substantial.
Cloud remains preferable only if power cost exceeds $\approx\$0.90$/kWh
or utilization falls below $\approx$20\%—conditions rare in well-run
enterprise data centers.

% ─────────────────────────────────────────────────────────────────────────────
\section{Implementation}
\label{sec:implementation}

InfraCalc is implemented as a single-file client-side web application
(HTML/CSS/JavaScript, $\approx$4,000 lines) with no server-side components.
All computations execute in the browser, enabling deployment as a static file
with zero backend infrastructure.

\paragraph{AI assistant.}
An embedded chatbot powered by the Cerebras API (model: \texttt{gpt-oss-120b})
explains calculator results to users. The assistant is context-injected on
each turn with pre-computed walkthroughs derived from the live calculator
state, ensuring it presents only verified numbers and never performs arithmetic
independently. Temperature is set to 0.3 to minimize creative deviation from
the injected context. A number validator post-processes model responses and
flags figures not present in the injected walkthrough.

\paragraph{Formula reference.}
A companion document (\texttt{InfraCalc\_Formula\_Guide.html}) provides
complete step-by-step derivations and worked examples for all 27 formulas,
accessible via the application's Formulas \& Assumptions tab and hyperlinked
from chatbot responses.

\paragraph{Reproducibility.}
The complete source is available at
\url{https://github.com/jaykolla/enterprise-ai-infracalc} under the MIT
license. All formulas are implemented directly in the open-source JavaScript
and match the derivations in this paper exactly.

% ─────────────────────────────────────────────────────────────────────────────
\section{Limitations and Future Work}
\label{sec:limitations}

\paragraph{Efficiency factor uncertainty.}
The throughput model uses a fixed efficiency factor $\eta{=}0.35$. Real
efficiency varies with batch size, model architecture, and inference runtime.
A natural extension is a lookup table of empirically measured $\eta$ values
per GPU, model, and runtime combination from MLCommons submissions.

\paragraph{Context management strategies.}
Our model assumes the full conversation context is retained in KV cache.
Production systems often apply sliding window attention, context summarization,
or prefix caching~\cite{zheng2023efficiently} that reduce effective context
below $T_{\text{eff}} = D \cdot (T_{\text{in}} + T_{\text{out}})$.
Modeling these strategies would tighten KV cache estimates for systems
that implement them.

\paragraph{Training workloads.}
InfraCalc focuses exclusively on inference serving. Fine-tuning and
full training workloads have different memory profiles (optimizer states,
gradient checkpointing) and are not currently modeled.

\paragraph{Multi-node networking.}
The model does not account for the communication overhead of multi-node
tensor parallelism, which affects throughput when the GPU count exceeds
a single server.

% ─────────────────────────────────────────────────────────────────────────────
\section{Conclusion}
\label{sec:conclusion}

We have presented InfraCalc, an end-to-end analytical framework for enterprise
LLM infrastructure planning that unifies GPU memory sizing, throughput and
latency estimation, TCO comparison, and unit economic analysis in a single
coherent model.

Our primary methodological contribution is the multi-turn conversation depth
multiplier (Equation~\ref{eq:kv_multiturn}), which corrects a systematic
sizing error in current practice. For a representative enterprise chatbot,
ignoring conversation depth produces a 10$\times$ underestimate of KV cache
and a 50\% underestimate of GPU count, leading to production failures.

The TCO model and normalized unit economic metrics—\$/GPU-hour and
\$/1M output tokens—provide a basis for rigorous on-premises versus cloud
decisions. Validation against MLCommons Inference v4.1 benchmarks confirms
that the throughput model is conservatively accurate within 8--14\% across
five GPU platforms.

InfraCalc is released as open-source software at
\url{https://github.com/jaykolla/enterprise-ai-infracalc} and we invite the
community to extend it with additional hardware platforms, model families,
and deployment scenarios.

% ─────────────────────────────────────────────────────────────────────────────
\section*{Acknowledgments}
The author thanks the open-source communities behind vLLM, MLCommons, and
the Llama model family. GPU hardware specifications referenced in this work
are sourced exclusively from publicly available vendor product pages and
datasheets.

% ─────────────────────────────────────────────────────────────────────────────
\bibliographystyle{plainnat}
\bibliography{references}

\end{document}
